\documentclass[11pt,a4paper]{article}
\usepackage[T1]{fontenc}
\usepackage{lmodern,amsmath,amssymb,amsthm,microtype}
\usepackage[margin=25mm]{geometry}
\usepackage[hidelinks]{hyperref}
\newtheorem{theorem}{Theorem}
\newtheorem{lemma}{Lemma}
\newtheorem{proposition}{Proposition}
\newcommand{\sqcupdot}{\bigsqcup}
\title{Domination polynomials of clique joins:\protect\\two-blade ultra-log-concavity and three-blade counterexamples}
\author{Aurantium}
\date{14 October 2026}
\begin{document}
\maketitle
\begin{abstract}
For all positive integers $r,a,b$, the domination polynomial of
$K_r\vee(K_a\sqcup K_b)$ is ultra-log-concave of order $N=r+a+b$,
and hence log-concave and unimodal. The proof uses an exact second
difference for normalized binomial coefficients, with the zero
constant coefficient and all truncation endpoints treated explicitly.
For three blades, ordinary log-concavity can fail:
$K_1\vee(K_4\sqcup K_5\sqcup K_7)$ has first coefficients
$(d_1,d_2,d_3)=(1,16,260)$. When equal blade sizes are allowed,
$K_1\vee(K_q\sqcup K_q\sqcup K_q)$ fails log-concavity for every
integer $q\ge5$. These results address parts of Rather's question
following Theorem~2.16; the seventeen-vertex example is itself
unimodal, so it is not a counterexample to general unimodality.
\end{abstract}

\section{Graphs, coefficient properties and the source question}
All graphs are finite, simple and undirected. A set $S$ of vertices
dominates a graph if every vertex outside $S$ has a neighbor in $S$.
Write
\[
 D_G(x)=\sum_{k=0}^{N}d_kx^k,
 \qquad d_k=\#\{S\subseteq V(G):|S|=k,\ S\text{ dominates }G\},
\]
where $N=|V(G)|$. Every graph considered here is nonempty, so $d_0=0$.

For positive integers $r,t,a_1,\ldots,a_t$, consider
\begin{equation}\label{eq:graph}
 G=K_r\vee\left(\sqcupdot_{i=1}^{t}K_{a_i}\right),
 \qquad M=\sum_{i=1}^{t}a_i,\qquad N=r+M.
\end{equation}
The join adds every edge from the central clique $K_r$ to each
blade; distinct blades have no edges between them. Blade orders may
coincide in this definition. The explicit counterexample below has
distinct blade orders $4,5,7$.

A nonnegative sequence $(d_k)_{k=0}^{N}$ is log-concave if
$d_k^2\ge d_{k-1}d_{k+1}$ for $1\le k\le N-1$. It is
\emph{ultra-log-concave of order $N$} if
\[
 u_k=\frac{d_k}{\binom Nk},\qquad 0\le k\le N,
\]
is log-concave and has no internal zeros. A sequence is unimodal if
it is nondecreasing up to some index and nonincreasing thereafter;
equal adjacent entries are allowed.

Rather asks about domination zeros and their limits, log-concavity
and unimodality for this clique-join family immediately after
Theorem~2.16 in \cite[p.~286]{R}. This is an open question in the
source, not an affirmative conjecture that all these properties
hold. We prove the full positive two-blade statement and give
three-blade failures of log-concavity. We do not resolve all parts
of that question.

\section{Direct counting and the polynomial formula}
\begin{lemma}\label{lem:count}
For the graph in \eqref{eq:graph},
\begin{equation}\label{eq:general}
 D_G(x)=\bigl((1+x)^r-1\bigr)(1+x)^M
              +\prod_{i=1}^{t}\bigl((1+x)^{a_i}-1\bigr).
\end{equation}
\end{lemma}
\begin{proof}
Partition dominating sets by whether they meet the central clique.
If they do, any nonempty subset of the central clique and any subset
of the other $M$ vertices dominate the whole graph: a selected central
vertex is adjacent to every other vertex. These sets contribute
$((1+x)^r-1)(1+x)^M$.

If no central vertex is selected, every blade must contribute a
nonempty subset, because vertices in a blade have no neighbors in
any other blade. Conversely, nonempty selections from every blade
dominate their own cliques and all central vertices. Here $t\ge1$,
so there is a selected neighbor for the central vertices.
The contribution is the product in \eqref{eq:general}.
The two classes are disjoint and exhaustive, proving the formula.
\end{proof}

Formula \eqref{eq:general} agrees with Theorem~2.16 of \cite{R}:
its first two terms combine as
\[
 \bigl((1+x)^r-1\bigr)\bigl((1+x)^M-1\bigr)
       +\bigl((1+x)^r-1\bigr)
 =\bigl((1+x)^r-1\bigr)(1+x)^M.
\]
The printed display in Corollary~2.17 contains an additional
$((1+x)^r-1)$ after this combined term, when its central-clique
parameter is renamed $r$. That display is not used here. All
counts follow the graph definition and Lemma~\ref{lem:count},
consistent with Theorem~2.16. This comparison does not assert an
author-issued correction to the source.

With two blades of orders $a,b$, expanding \eqref{eq:general} gives
\begin{equation}\label{eq:two}
 D_G(x)=(1+x)^N-(1+x)^a-(1+x)^b+1,
 \qquad N=r+a+b.
\end{equation}
Thus
\begin{equation}\label{eq:coeff}
 d_0=0,\qquad
 d_k=\binom Nk-\binom ak-\binom bk\quad(1\le k\le N).
\end{equation}
We use $\binom ak=0$ for integers $k>a$, and likewise for $b$.
For each $1\le k\le N$, fixing any one central vertex and choosing
$k-1$ of the other vertices yields dominating sets, so
\begin{equation}\label{eq:positive}
 d_k\ge\binom{N-1}{k-1}>0.
\end{equation}

\section{Ultra-log-concavity for every two-blade join}
\begin{theorem}\label{thm:ulc}
For every triple of positive integers $r,a,b$, the domination
coefficient sequence of $K_r\vee(K_a\sqcup K_b)$ is ultra-log-concave
of order $N=r+a+b$. It is therefore log-concave and unimodal.
\end{theorem}
\begin{proof}
Because $r,a,b\ge1$, we have $N\ge3$ and
$N-a=r+b\ge2$, $N-b=r+a\ge2$. Define, for $0\le k\le N$,
\[
 A_k=\frac{\binom ak}{\binom Nk},\qquad
 B_k=\frac{\binom bk}{\binom Nk}.
\]
The denominator is strictly positive throughout this index range.
By \eqref{eq:coeff}, the normalized coefficients satisfy
\begin{equation}\label{eq:normalized}
 u_0=0,\qquad u_k=1-A_k-B_k\quad(1\le k\le N).
\end{equation}
In particular, the latter expression is not to be substituted at
$k=0$, where $1-A_0-B_0=-1$ rather than $u_0$.

We claim that, for every $0\le k\le N-2$,
\begin{equation}\label{eq:second}
 A_k-2A_{k+1}+A_{k+2}
 =A_k\frac{(N-a)(N-a-1)}{(N-k)(N-k-1)}\ge0.
\end{equation}
For $k\le a$, the two needed ratios are
\[
 A_{k+1}=A_k\frac{a-k}{N-k},\qquad
 A_{k+2}=A_k\frac{(a-k)(a-k-1)}{(N-k)(N-k-1)}.
\]
These equalities include the zero-factor cases $k=a-1$ and $k=a$;
they do not divide by a vanishing numerator or by $A_{k+1}$.
All denominators are positive because $k\le N-2$.
Writing $s=N-k$ and $d=N-a$, their common numerator simplifies by
\[
 s(s-1)-2(s-d)(s-1)+(s-d)(s-d-1)=d(d-1),
\]
which proves \eqref{eq:second} in this range. For $k>a$, all three
terms $A_k,A_{k+1},A_{k+2}$ are zero by the binomial convention,
and the same identity holds. This covers every boundary case.
Replacing $a$ by $b$ proves the identical statement for $B$.

For $2\le k\le N-1$, all three indices $k-1,k,k+1$ are positive,
so \eqref{eq:normalized} and the second-difference identities at
index $k-1$ give
\[
 u_{k-1}-2u_k+u_{k+1}\le0.
\]
The entries $u_1,\ldots,u_N$ are positive by \eqref{eq:positive}.
The arithmetic-geometric mean inequality therefore gives
\begin{equation}\label{eq:AGM}
 u_k\ge\frac{u_{k-1}+u_{k+1}}2
       \ge\sqrt{u_{k-1}u_{k+1}}
 \qquad(2\le k\le N-1).
\end{equation}
For the remaining index $k=1$, the required log-concavity inequality
is automatic because $u_0=0$:
$u_1^2\ge0=u_0u_2$. Thus $(u_k)_{k=0}^{N}$ satisfies every required
log-concavity inequality, and its positive support is $1,\ldots,N$.
This proves ultra-log-concavity of order $N$.

Finally the binomial row is positive and log-concave, since
\[
 \frac{\binom Nk^2}{\binom N{k-1}\binom N{k+1}}
 =\frac{(k+1)(N-k+1)}{k(N-k)}\ge1\qquad(1\le k\le N-1).
\]
Multiplying
its log-concavity inequality by that of $u$ proves
$d_k^2\ge d_{k-1}d_{k+1}$ for $1\le k\le N-1$.
Since $d_1,\ldots,d_N>0$, their successive ratios are nonincreasing
and can pass from at least one to at most one only once. Hence
the sequence is unimodal, with any equality plateau allowed;
the leading $d_0=0$ does not change that conclusion.
\end{proof}

At the smallest parameters $r=a=b=1$, formula \eqref{eq:two} gives
$(d_0,d_1,d_2,d_3)=(0,1,3,1)$ and
$(u_0,u_1,u_2,u_3)=(0,1/3,1,1)$. This illustrates the need to
treat $k=1$ through $u_0=0$: the concavity argument concerns only
$u_1,\ldots,u_N$, not the entire normalized sequence starting at zero.

\section{Three-blade failures of ordinary log-concavity}
\begin{proposition}\label{prop:17}
The graph $G=K_1\vee(K_4\sqcup K_5\sqcup K_7)$ has $17$ vertices
and a domination polynomial that is not log-concave.
\end{proposition}
\begin{proof}
A dominating set of size one or two must contain the central vertex:
otherwise each of the three blades would need a selected vertex.
Thus $d_1=1$ and $d_2=16$. A dominating set of size three either
contains the central vertex and two arbitrary other vertices, or
contains one vertex from each blade. These disjoint cases give
\[
 d_3=\binom{16}{2}+4\cdot5\cdot7=120+140=260.
\]
Consequently $d_2^2-d_1d_3=256-260=-4<0$, which violates
log-concavity at index two.
\end{proof}

This counterexample has distinct blade orders, so it does not
depend on an interpretation allowing equal orders in the source's
wording. It is also \emph{unimodal}. For completeness, its full
coefficient sequence, equivalently obtained by expanding
\[
 x(1+x)^{16}+
 \bigl((1+x)^4-1\bigr)\bigl((1+x)^5-1\bigr)\bigl((1+x)^7-1\bigr),
\]
is
\begin{align*}
 (d_0,\ldots,d_{17})={}&(0,1,16,260,1470,4830,10913,18291,23641,\\
                     &24034,19371,12363,6187,2380,680,136,17,1).
\end{align*}
It increases up to $d_9=24034$ and then decreases. These full
coefficients also occur in the earlier exact graph enumeration;
the direct three-coefficient proof above does not require that
enumeration.

\begin{proposition}\label{prop:equal}
When equal blade orders are allowed, every graph
$K_1\vee(K_q\sqcup K_q\sqcup K_q)$ with integer $q\ge5$ has a
domination polynomial that is not log-concave.
\end{proposition}
\begin{proof}
The same partition of sets of sizes one, two and three gives
\[
 d_1=1,\qquad d_2=3q,\qquad d_3=\binom{3q}{2}+q^3.
\]
Their log-concavity difference is exactly
\begin{equation}\label{eq:equal-failure}
 d_2^2-d_1d_3=-\frac q2(2q^2-9q-3).
\end{equation}
The factor $2q^2-9q-3$ is $2$ at $q=5$, and its increment when
$q$ increases to $q+1$ is $4q-7>0$ for $q\ge5$.
Thus \eqref{eq:equal-failure} is strictly negative for every
integer $q\ge5$. At $q=5$ the first three coefficients are
$(1,15,230)$, giving difference $-5$.
\end{proof}

The equal-blade family is stated for the graph definition
\eqref{eq:graph}, which explicitly permits repeated orders. Its
negative conclusion concerns log-concavity; no claim that this
family fails unimodality is made.

\section{Evidence and limits of the conclusions}
All parameter-wide results above follow from the written counting
and algebraic proofs. Earlier exact checks covered all $1728$
ordered triples $r,a,b\in\{1,\ldots,12\}$, including $63936$
second-difference equalities and $31968$ ultra-log-concavity
inequalities. Separate earlier vertex-subset enumeration checked
four small two-blade graphs and the full coefficient vectors of
the two displayed three-blade examples. These are historical
finite checks, not new executions for this note or substitutes
for the general proof. Their precise execution scope is recorded
in the accompanying reproducibility note; no program is needed
for the mathematical argument and no larger suite is bundled.

Theorem~\ref{thm:ulc} covers exactly positive $r,a,b$ and two blades.
The three-blade counterexamples rule out universal log-concavity,
but do not resolve unimodality for arbitrary $t\ge3$ or the
source's general questions about zeros and their limits. No claim
of smallest counterexample order, first discovery or worldwide
priority is made. Prepared with AI assistance; no proof-assistant
certification or independent human peer review is asserted.

\begin{thebibliography}{1}
\raggedright
\bibitem{R}
Bilal Ahmad Rather, \emph{On domination polynomials of some graphs},
Journal of Combinatorial Mathematics and Combinatorial Computing
\textbf{126} (2025), 279--289.
\href{https://doi.org/10.61091/jcmcc126-19}{doi:10.61091/jcmcc126-19}.
Theorem~2.16, its following open question, and Corollary~2.17,
printed p.~286.
\end{thebibliography}
\end{document}
